\section*{Hoofdstuk \ref{ch:b2:catalytic-cycles} --- Organometaalkatalyse: elementaire stappen en cycli}

\begin{solution}{exo:b2:catalytic-cycles:1}
Ervoor: Ir(I), $d^8$, $8 + 4 \times 2 = 16$ elektronen. Erna: Ir(III), $d^6$, $6 + 6 \times 2 = 18$
elektronen, octaëdrisch.
\end{solution}

\begin{solution}{exo:b2:catalytic-cycles:2}
De eerste is een \omterm{def:b2:catalytic-cycles:oxidative-addition}{reductieve eliminatie} (Pt(IV) naar Pt(II), ethaan gevormd uit twee methylgroepen).
De tweede is een \omterm{def:b2:catalytic-cycles:ligand-exchange}{liganduitwisseling} (één fosfine vervangen door etheen) gevolgd door een
\omterm{def:b2:catalytic-cycles:insertion}{migratorische insertie} van etheen in de binding Pd--Ph.
\end{solution}

\begin{solution}{exo:b2:catalytic-cycles:3}
TON $= 15.0/0.020 = 750$; TOF $= 750/3.0 = \qty{250}{h^{-1}}$.
\end{solution}

\begin{solution}{exo:b2:catalytic-cycles:4}
\begin{align*}
  &\mathrm{PdL_2 + ArBr \to ArPdBrL_2}\\
  &\mathrm{ArPdBrL_2 + Ar'B(OH)_3^- \to ArPdAr'L_2 + Br^- + B(OH)_3}\\
  &\mathrm{ArPdAr'L_2 \to Ar{-}Ar' + PdL_2}
\end{align*}
Som: $\mathrm{ArBr + Ar'B(OH)_3^- \to Ar{-}Ar' + Br^- + B(OH)_3}$; elk palladiumdeeltje
valt weg.
\end{solution}

\begin{solution}{exo:b2:catalytic-cycles:5}
Ethyl en isopropyl (die hebben $\beta$-waterstofatomen). Methyl heeft geen $\beta$-koolstofatoom; het $\beta$-koolstofatoom
van neopentyl draagt geen waterstof; het $\beta$-koolstofatoom van benzyl is een \omterm{def:b2:huckel:aromatic}{aromatisch} ringkoolstofatoom zonder
waterstof.
\end{solution}

\begin{solution}{exo:b2:catalytic-cycles:6}
Methyl-(\textit{E})-3-fenylpropenoaat, \ce{Ph-CH=CH-COOCH3}, met de fenylgroep op het
eindstandige koolstofatoom, trans-product.
\end{solution}

\begin{solution}{exo:b2:catalytic-cycles:7}
Butanal \ce{CH3CH2CH2CHO} (lineair: het metaal addeert aan het eindstandige koolstofatoom, het hydride aan
het binnenste) en 2-methylpropanal \ce{(CH3)2CHCHO} (vertakt: het metaal op het binnenste
koolstofatoom).
\end{solution}

\begin{solution}{exo:b2:catalytic-cycles:8}
\ce{[RhCl(PPh3)3]} heeft 16 elektronen: \ce{H2} adderen (+2) zou 18 geven met zes
coördinatieposities, te druk bezet met drie omvangrijke fosfinen. \ce{[RhCl(PPh3)2]} heeft er 14
en een vrije plaats: \omterm{def:b2:catalytic-cycles:oxidative-addition}{oxidatieve additie} geeft gemakkelijk een complex met 16 elektronen.
\end{solution}

\begin{solution}{exo:b2:catalytic-cycles:9}
Het boorzuur is een zwak nucleofiel; de base voegt hydroxide (of alkoxide) toe aan boor,
wat een boraat \ce{ArB(OH)3-} geeft, waarvan de arylgroep elektronenrijker is en naar
palladium overgaat.
\end{solution}

\begin{solution}{exo:b2:catalytic-cycles:10}
Een complex met 20 elektronen is niet aannemelijk. Het complex moet eerst een ligand verliezen (14 e), daarna
\ce{H2} adderen (16 e), het alkeen binden (18 e), het in een binding M--H inserteren (16 e) en het
alkaan elimineren (14 e), voordat het het ligand terugneemt of opnieuw begint.
\end{solution}

\begin{solution}{exo:b2:catalytic-cycles:11}
Beginsnelheid $n_{\max}/\tau = \qty{0.25}{mmol/min}$; begin-TOF $= 0.25/0.010 =
\qty{25}{min^{-1}} = \qty{1.5e3}{h^{-1}}$. Uiteindelijk TON $= 10.0/0.010 = 1.0 \times 10^3$.
\end{solution}

\begin{solution}{exo:b2:catalytic-cycles:12}
4-Broomanisool met fenylboorzuur (of broombenzeen met 4-methoxyfenylboorzuur),
een palladium(0)-fosfinekatalysator en een base. Cyclus: \omterm{def:b2:catalytic-cycles:oxidative-addition}{oxidatieve additie} van het
arylbromide, \omterm{def:b2:catalytic-cycles:transmetalation}{transmetallering} vanuit het boraat, \omterm{def:b2:catalytic-cycles:oxidative-addition}{reductieve eliminatie} van
4-methoxybifenyl.
\end{solution}

\begin{solution}{pb:b2:catalytic-cycles:1}
\textbf{1.} Pd(II), $d^8$.
\textbf{2.} Pd(0), $d^{10}$, $10 + 2 \times 2 = 14$ elektronen.
\textbf{3.} Met 14 elektronen kan het nog twee paren opnemen: het heeft vrije coördinatieplaatsen.
\textbf{4.} Palladium(0)-fosfinecomplexen worden door lucht geoxideerd, en fosfinen worden
geoxideerd tot fosfineoxiden.
\textbf{5.} Het reduceert palladium(II) tot palladium(0) (en wordt daarbij zelf geoxideerd) en bindt het
palladium(0), zodat het in oplossing blijft.
\textbf{6.} De \omterm{def:b2:catalytic-cycles:cycle}{prekatalysator}.
\textbf{7.} \ce{Pd(PPh3)2 + CH3C6H4Br -> [Pd(C6H4CH3)Br(PPh3)2]}: Pd(II), $d^8$, $8 + 4 \times 2 =
16$ elektronen.
\textbf{8.} Het zet het om in het boraat \ce{PhB(OH)3-}, dat zijn fenylgroep overdraagt.
\textbf{9.} \ce{[Pd(Ar)Br(PPh3)2] + PhB(OH)3- -> [Pd(Ar)(Ph)(PPh3)2] + Br- + B(OH)3}.
\textbf{10.} \omterm{def:b2:catalytic-cycles:oxidative-addition}{Reductieve eliminatie} verbindt twee cis staande liganden; trans staande arylgroepen moeten eerst isomeriseren.
\textbf{11.} \ce{[Pd(Ar)(Ph)(PPh3)2] -> Ar-Ph + Pd(PPh3)2}: palladium(0), 14 elektronen, klaar
voor een nieuwe omloop.
\textbf{12.} \ce{CH3C6H4Br + PhB(OH)3- -> CH3C6H4-Ph + Br- + B(OH)3}.
\textbf{13.} 4-Broomtolueen (\qty{10.0}{mmol}, tegenover \qty{12.0}{mmol} boorzuur).
\textbf{14.} $1.51/168.24 = \qty{8.98e-3}{mol}$, \qty{8.98}{mmol}.
\textbf{15.} $8.98/10.0 = 89.8~\%$.
\textbf{16.} $0.050/10.0 = 0.50$~mol~\%.
\textbf{17.} $8.98/0.050 = 180$.
\textbf{18.} $180/4.0 = \qty{45}{h^{-1}}$.
\textbf{19.} Twee fenylgroepen die naar hetzelfde palladiumatoom overgedragen worden (homokoppeling van het
boorzuur, begunstigd door sporen zuurstof), en daarna samen geëlimineerd worden.
\textbf{20.} De arylgroepen hebben geen $\beta$-waterstofatoom dat het metaal kan bereiken.
\textbf{21.} De \omterm{def:b2:catalytic-cycles:oxidative-addition}{oxidatieve additie}: de binding \ce{C-Cl} is sterker dan \ce{C-Br}.
\textbf{22.} Het eindigt als palladium(0)-deeltjes of -zouten in de residuen; het is duur en
giftig, en wordt teruggewonnen.
\textbf{23.} Een deel van het boorzuur gaat verloren aan nevenreacties (homokoppeling, verlies van boor
van de ring); de overmaat houdt het bromide beperkend.
\textbf{24.} Het palladium heeft $\boldsymbol{\approx \num{1.8e2}}$ omlopen gemaakt.
\end{solution}
